% generated by roam.py handbook export -- не редактировать, править заметки

\chapter{Вспомогательные утверждения}\label{aux}
\InputIfFileExists{text/aux-intro}{}{}

\section{К лекции 1}\label{aux:1}

\InputIfFileExists{text/1B626AA4-35C8-4C7C-9D3D-D8C26D98F40C-before}{}{}
% generated by roam.py handbook export -- не редактировать, править заметки
% 20260916214817-векторное_пространство.org (body 1522a4bcd559)

\begin{definition}[{Векторное пространство}]\label[definition]{note:1B626AA4-35C8-4C7C-9D3D-D8C26D98F40C}
  Пусть $K$~--- поле (в курсе $K = \R$, реже $\C$). Множество $V$ с операциями сложения $V \times V \rightarrow V$ и умножения на скаляр $K \times V \rightarrow V$ называется \textbf{векторным} (линейным) пространством над $K$, если $\forall x, y, z \in V\ \forall \alpha, \beta \in K$:

  \begin{enumerate}
    \item $(x + y) + z = x + (y + z)$
    \item $x + y = y + x$
    \item $\exists\, 0 \in V:\ x + 0 = x$
    \item $\forall x\ \exists\, (-x) \in V:\ x + (-x) = 0$
    \item $\alpha (\beta x) = (\alpha \beta) x$
    \item $1 \cdot x = x$
    \item $\alpha (x + y) = \alpha x + \alpha y$
    \item $(\alpha + \beta) x = \alpha x + \beta x$
  \end{enumerate}

  Примеры: $\R^n$; матрицы $\R^{m \times n}$; непрерывные функции $C[a;b]$; многочлены степени не выше $k$.
\end{definition}

\begin{remark}
  В курсе почти всегда $V = \R^n$. Для \hyperref[note:316B5BFA-8B9F-4267-9086-D6593AAB5F9B]{выпуклых множеств} важно, что поле вещественное: определение использует $\lambda \in [0;1]$ и комбинации $\lambda x + (1 - \lambda) y$.
\end{remark}

\begin{remark}
  Дополнительные структуры на $V$: \hyperref[note:F2942E98-6115-4600-AA3F-E0F687D9C732]{скалярное произведение} (и порожденная им норма $\norm{x} = \sqrt{\left\langle x; x \right\rangle}$), полнота по этой норме~--- \hyperref[note:5D688904-8445-47B5-9D39-8BCC39F9C354]{гильбертово пространство}.
\end{remark}
  % Векторное пространство -- пространство $V$ в определениях раздела 1.1

\InputIfFileExists{text/3819FAC4-B7DB-4181-A40B-8FBB5EC85DC1-before}{}{}
% generated by roam.py handbook export -- не редактировать, править заметки
% 20260917205407-борелевская_оболочка.org (body af22c31c0e94)

\begin{definition}[{Борелевская оболочка}]\label[definition]{note:3819FAC4-B7DB-4181-A40B-8FBB5EC85DC1}
  Пусть $X$~--- множество, $\mathcal{E}$~--- система его подмножеств. \textbf{Борелевская оболочка} $\mathcal{B}(\mathcal{E})$~--- наименьшая $\sigma$-алгебра подмножеств $X$, содержащая $\mathcal{E}$:
  \[
  \mathcal{B}(\mathcal{E}) = \bigcap \left\{ \mathcal{A}\ \middle|\ \mathcal{A} \text{ -- } \sigma\text{-алгебра подмножеств } X,\ \mathcal{E} \subset \mathcal{A} \right\}
  \]

  Напоминание: система $\mathcal{A}$ подмножеств $X$~--- $\sigma$-алгебра, если

  \begin{enumerate}
    \item $X \in \mathcal{A}$;
    \item $A \in \mathcal{A} \Rightarrow X \setminus A \in \mathcal{A}$;
    \item $A_1, A_2, \dots \in \mathcal{A} \Rightarrow \bigcup\limits_{k=1}^{\infty} A_k \in \mathcal{A}$.
  \end{enumerate}
\end{definition}

\begin{remark}
  Главный пример~--- \emph{борелевская} $\sigma$-алгебра $\mathfrak{B}(\R^n)$: борелевская оболочка системы всех открытых множеств $\R^n$.
\end{remark}

\begin{remark}
  Та же схема «пересечь все объекты с нужным свойством, содержащие $M$» дает \hyperref[note:1AC657F9-2A3D-47D1-A71D-5B813D961260]{выпуклую оболочку}, линейную оболочку, замыкание множества. Она работает, когда (1) хотя бы один такой объект есть и (2) свойство сохраняется при любых пересечениях~--- как у $\sigma$-алгебр (ниже) и у \hyperref[note:08CF7113-4709-4C4F-A0E7-4A2ED595BD7F]{выпуклых множеств}.
\end{remark}

\begin{proof}
  Семейство в пересечении непусто: $2^X$ (все подмножества $X$)~--- $\sigma$-алгебра, содержащая $\mathcal{E}$.

  Пересечение $\mathcal{B}$ любого семейства $\sigma$-алгебр $\mathcal{A}_i$~--- снова $\sigma$-алгебра: каждая аксиома выполнена в каждой $\mathcal{A}_i$, поэтому и в пересечении. Например, если $A_1, A_2, \dots \in \mathcal{B}$, то все $A_k$ лежат в каждой $\mathcal{A}_i$, значит $\bigcup A_k \in \mathcal{A}_i$ для всех $i$, то есть $\bigcup A_k \in \mathcal{B}$.

  Минимальность: $\mathcal{E} \subset \mathcal{B}(\mathcal{E})$, и любая $\sigma$-алгебра $\mathcal{A} \supset \mathcal{E}$ участвует в пересечении, поэтому $\mathcal{B}(\mathcal{E}) \subset \mathcal{A}$.
\end{proof}
  % Борелевская оболочка -- та же конструкция «пересечение всех объектов, содержащих $M$», что и в опр. 1.2

\InputIfFileExists{text/14FE7B45-F304-4C15-8086-3535589FA04A-before}{}{}
% generated by roam.py handbook export -- не редактировать, править заметки
% 20260916214241-сумма_минковского.org (body fa864640c037)

\begin{definition}[{Сумма Минковского}]\label[definition]{note:14FE7B45-F304-4C15-8086-3535589FA04A}
  Сумма Минковского для двух подмножеств векторного пространства $V$, пусть это будут $A,B \subset V$~--- это множество, состоящее из всевозможных сумм элементов этих подмножеств.

  $A + B = \left\{ a + b \mid a \in A,\ b \in B \right\}$
\end{definition}
  % Сумма Минковского -- для свойства 1.1.3

\InputIfFileExists{text/5D688904-8445-47B5-9D39-8BCC39F9C354-before}{}{}
% generated by roam.py handbook export -- не редактировать, править заметки
% 20260916132456-гильбертово_пространство.org (body 2ffefebdf65e)

\begin{definition}[{Гильбертово пространство}]\label[definition]{note:5D688904-8445-47B5-9D39-8BCC39F9C354}
  Это векторное пространство, на котором задано \hyperref[note:F2942E98-6115-4600-AA3F-E0F687D9C732]{скалярное произведение}, полное относительно нормы $\norm{x} = \sqrt{\left\langle x; x \right\rangle}$
\end{definition}
  % Гильбертово пространство

\InputIfFileExists{text/F2942E98-6115-4600-AA3F-E0F687D9C732-before}{}{}
% generated by roam.py handbook export -- не редактировать, править заметки
% 20260916132628-скалярное_произведение.org (body c6c700890596)

\begin{definition}[{Скалярное произведение}]\label[definition]{note:F2942E98-6115-4600-AA3F-E0F687D9C732}
  Скалярное произведение: функция 2 аргументов, удовлетворяющая следующим свойствам.

  Пусть $X$~--- векторное пространство над $\mathbb{R}$ или $\mathbb{C}$.
  Тогда
  $\left\langle \cdot, \cdot \right\rangle:\ X \times X \rightarrow \mathbb{R} \text{ или } \mathbb{C}$~--- скалярное произведение, если:

  \begin{enumerate}
    \item $\left\langle x; x \right\rangle \ge 0\ \forall x,\ \left\langle x; x \right\rangle = 0 \Leftrightarrow x = \overline{0}$
    \item $\left\langle x+y; z \right\rangle = \left\langle x; z \right\rangle + \left\langle y; z \right\rangle\ \forall x,y,z \in X$
    \item $\left\langle c \cdot x; y \right\rangle = c \cdot \left\langle x; y \right\rangle, \forall x, y \in X, \forall c \in \mathbb{R},\ \mathbb{C}$
    \item $\left\langle x; y \right\rangle = \overline{\left\langle y ;x \right\rangle}\ \forall x, y \in X$
  \end{enumerate}
\end{definition}
  % Скалярное произведение -- для свойства 1.1.2

\InputIfFileExists{text/B9EF495E-CDBE-4E80-B401-8AFEAB0DCCAC-before}{}{}
% generated by roam.py handbook export -- не редактировать, править заметки
% 20260918190111-аффинные_преобразования.org (body 09d775c9f42c)

\begin{definition}[{Аффинные преобразования}]\label[definition]{note:B9EF495E-CDBE-4E80-B401-8AFEAB0DCCAC}
  Пусть $V, W$~--- векторные пространства. Отображение $\mathcal{A}: V \rightarrow W$~--- \emph{аффинное}, если $\mathcal{A}(x) = Ax + b$, где $A: V \rightarrow W$~--- линейный оператор, $b \in W$.
\end{definition}

\begin{proposition}[{Свойство: сохраняет выпуклые комбинации}]\label[proposition]{note:B9EF495E-CDBE-4E80-B401-8AFEAB0DCCAC:1}
  Для любых $x, y \in V$ и $\lambda \in \R$
  \[
  \mathcal{A}(\lambda x + (1-\lambda) y) = \lambda \mathcal{A}(x) + (1-\lambda) \mathcal{A}(y)
  \]
  Следствие: образ и прообраз \hyperref[note:316B5BFA-8B9F-4267-9086-D6593AAB5F9B]{выпуклого множества} при аффинном отображении~--- выпуклые.
\end{proposition}

\begin{proof}
  По линейности $A$: $\mathcal{A}(\lambda x + (1-\lambda) y) = \lambda Ax + (1-\lambda) Ay + b$. Так как $\lambda + (1-\lambda) = 1$, то $b = \lambda b + (1-\lambda) b$, и правая часть равна $\lambda (Ax + b) + (1-\lambda)(Ay + b)$.
\end{proof}
  % Аффинные преобразования -- для утверждения 1.2, п. 1

\InputIfFileExists{text/1756907B-1FF2-4F62-8553-759B8509281C-before}{}{}
% generated by roam.py handbook export -- не редактировать, править заметки
% 20260918211206-градиент.org (body 60481c528ed3)

\begin{definition}[{Градиент}]\label[definition]{note:1756907B-1FF2-4F62-8553-759B8509281C}
  Пусть $E \subset \R^n$~--- открытое множество, $f: E \rightarrow \R$~--- \hyperref[note:957FA09D-8883-40A5-B4C7-237D9F815A42]{дифференцируемая} в точке $x \in E$ функция.
  \textbf{Градиент} $f$ в точке $x$~--- вектор-столбец из частных производных:
  \[
  \nabla f(x) = \left( \frac{\partial f}{\partial x_1}(x), \dots, \frac{\partial f}{\partial x_n}(x) \right)^T \in \R^n
  \]

  Производная $f^{\prime}(x)$~--- это строка $1 \times n$ (матрица линейного отображения $\R^n \rightarrow \R$), а градиент~--- она же, записанная столбцом: $\nabla f(x) = (f^{\prime}(x))^T$.
  Поэтому действие производной на вектор $h \in \R^n$~--- это \hyperref[note:F2942E98-6115-4600-AA3F-E0F687D9C732]{скалярное произведение} с градиентом:
  \[
  f^{\prime}(x) h = (\nabla f(x))^T h = \left\langle \nabla f(x); h \right\rangle
  \]
  и определение дифференцируемости переписывается как
  \[
  f(x + h) = f(x) + \left\langle \nabla f(x); h \right\rangle + o(\norm{h}),\ h \rightarrow 0
  \]
\end{definition}

\begin{proposition}[{Производная вдоль прямой}]\label[proposition]{note:1756907B-1FF2-4F62-8553-759B8509281C:1}
  Пусть $x \in E$, $h \in \R^n$ и $g(t) = f(x + th)$ (для тех $t$, при которых $x + th \in E$). Тогда
  \[
  g^{\prime}(t) = \left\langle \nabla f(x + th); h \right\rangle
  \]
\end{proposition}

\begin{proof}
  $g = f \circ l$, где $l(t) = x + th$, $l: \R \rightarrow \R^n$ и $l^{\prime}(t) = h$ (столбец $n \times 1$). По правилу дифференцирования композиции:
  \[
  g^{\prime}(t) = f^{\prime}(l(t)) l^{\prime}(t) = (\nabla f(x + th))^T h = \left\langle \nabla f(x + th); h \right\rangle
  \]
\end{proof}
  % Градиент

\InputIfFileExists{text/592AD816-1D80-48E6-8399-8524FBC1E8F3-before}{}{}
% generated by roam.py handbook export -- не редактировать, править заметки
% 20260920221509-матрица_гессе.org (body 7019915fa82d)

\begin{definition}[{Матрица Гессе}]\label[definition]{note:592AD816-1D80-48E6-8399-8524FBC1E8F3}
  Пусть $E \subset \R^n$~--- открытое множество, $f: E \rightarrow \R$~--- дважды непрерывно дифференцируемая функция.
  \textbf{Матрица Гессе} (\emph{гессиан}) $f$ в точке $x \in E$~--- матрица $n \times n$ из вторых частных производных:
  \[
  \nabla^2 f(x) = \left( \frac{\partial^2 f}{\partial x_i \partial x_j}(x) \right)_{i,j=1}^n
  \]

  Смешанные производные непрерывны, поэтому не зависят от порядка дифференцирования, и $\nabla^2 f(x)$~--- симметричная матрица.

  Матрица Гессе~--- это производная \hyperref[note:1756907B-1FF2-4F62-8553-759B8509281C]{градиента} как отображения $\nabla f: E \rightarrow \R^n$: ее $i$-я строка состоит из частных производных $i$-й координаты градиента $\frac{\partial f}{\partial x_i}$.

  Запись $\nabla^2 f(x) \ge 0$ означает \emph{неотрицательную определенность}: $h^T \nabla^2 f(x) h \ge 0\ \forall h \in \R^n$. Это нестрогий аналог \hyperref[note:771B93FA-654D-4233-9487-B4FDD356DDFE]{положительной определенности}, где требуется $h^T A h > 0\ \forall h \ne 0$.
\end{definition}

\begin{proposition}[{Вторая производная вдоль прямой}]\label[proposition]{note:592AD816-1D80-48E6-8399-8524FBC1E8F3:1}
  Пусть $x \in E$, $h \in \R^n$ и $\phi(t) = f(x + th)$ (для тех $t$, при которых $x + th \in E$). Тогда
  \[
  \phi^{\prime\prime}(t) = \left\langle \nabla^2 f(x + th) h; h \right\rangle = h^T \nabla^2 f(x + th) h
  \]
\end{proposition}

\begin{proof}
  Из заметки про \hyperref[note:1756907B-1FF2-4F62-8553-759B8509281C]{градиент} (производная вдоль прямой):
  \[
  \phi^{\prime}(t) = \left\langle \nabla f(x + th); h \right\rangle = \sum\limits_{i=1}^n \frac{\partial f}{\partial x_i}(x + th) h_i
  \]
  Каждое слагаемое~--- снова функция вида $u(x + th)$, где $u = \frac{\partial f}{\partial x_i}$, поэтому к нему применима та же формула:
  \[
  \left( \frac{\partial f}{\partial x_i}(x + th) \right)^{\prime}_t = \left\langle \nabla \frac{\partial f}{\partial x_i}(x + th); h \right\rangle = \sum\limits_{j=1}^n \frac{\partial^2 f}{\partial x_i \partial x_j}(x + th) h_j
  \]
  Подставим:
  \[
  \phi^{\prime\prime}(t) = \sum\limits_{i=1}^n \sum\limits_{j=1}^n \frac{\partial^2 f}{\partial x_i \partial x_j}(x + th) h_i h_j = h^T \nabla^2 f(x + th) h
  \]
\end{proof}
  % Матрица Гессе -- производные вдоль прямой в теоремах 1.2, 1.3

\InputIfFileExists{text/2F475EE4-2630-45FD-8DD2-8EA9D48059E4-before}{}{}
% generated by roam.py handbook export -- не редактировать, править заметки
% 20260918211205-дифференциальныи_критерии_выпуклости_функции_однои_переменнои.org (body abe91812a6bc)

\begin{theorem}[{Дифференциальный критерий выпуклости функции одной переменной}]\label[theorem]{note:2F475EE4-2630-45FD-8DD2-8EA9D48059E4}
  Пусть $I \subset \R$~--- промежуток, $g: I \rightarrow \R$~--- дифференцируемая функция.
  Тогда $g$~--- \hyperref[note:1CA14F2C-637F-41DE-9E97-72075D19E97F]{выпуклая} $\Leftrightarrow$ $\forall s,t \in I\ g(t) \ge g(s) + g^{\prime}(s)(t-s)$
\end{theorem}

\begin{remark}
  Геометрически: график выпуклой функции лежит не ниже любой своей касательной.
\end{remark}

\begin{proof}
  \proofpart{$\Rightarrow$}
  Пусть $g$~--- выпуклая, $s \ne t$ (при $s = t$ неравенство очевидно). Возьмем $\lambda \in (0;1]$, по выпуклости:
  \[
  g(s + \lambda(t-s)) = g((1-\lambda)s + \lambda t) \le (1-\lambda) g(s) + \lambda g(t)
  \]
  Перенесем $g(s)$ влево и поделим на $\lambda > 0$:
  \[
  \frac{g(s + \lambda(t-s)) - g(s)}{\lambda} \le g(t) - g(s)
  \]
  Домножим и поделим левую часть на $t-s$ и устремим $\lambda \rightarrow +0$:
  \[
  \frac{g(s + \lambda(t-s)) - g(s)}{\lambda (t-s)} (t-s) \rightarrow g^{\prime}(s)(t-s)
  \]
  Нестрогое неравенство сохраняется в пределе, значит $g^{\prime}(s)(t-s) \le g(t) - g(s)$.

  \proofpart{$\Leftarrow$}
  Возьмем $s,t \in I$, $\lambda \in [0;1]$ и точку $z = \lambda s + (1-\lambda) t \in I$. Запишем неравенство из условия в точке $z$ дважды:
  \[
  g(s) \ge g(z) + g^{\prime}(z)(s-z)
  \]
  \[
  g(t) \ge g(z) + g^{\prime}(z)(t-z)
  \]
  Умножим первое на $\lambda$, второе на $1-\lambda$ и сложим:
  \[
  \lambda g(s) + (1-\lambda) g(t) \ge g(z) + g^{\prime}(z)(\lambda s + (1-\lambda) t - z) = g(z)
  \]
  Это и есть определение выпуклости: $g(\lambda s + (1-\lambda)t) \le \lambda g(s) + (1-\lambda) g(t)$.
\end{proof}
  % Дифференциальный критерий выпуклости функции одной переменной

\InputIfFileExists{text/9CEE854D-ED37-4C08-82EB-E41B36C43EA6-before}{}{}
% generated by roam.py handbook export -- не редактировать, править заметки
% 20260920221510-критерии_выпуклости_функции_однои_переменнои_через_вторую_производную.org (body 02f42c324a24)

\begin{theorem}[{Критерий выпуклости функции одной переменной через вторую производную}]\label[theorem]{note:9CEE854D-ED37-4C08-82EB-E41B36C43EA6}
  Пусть $I \subset \R$~--- промежуток, $\phi: I \rightarrow \R$~--- дважды дифференцируемая функция.
  Тогда $\phi$~--- \hyperref[note:1CA14F2C-637F-41DE-9E97-72075D19E97F]{выпуклая} $\Leftrightarrow$ $\phi^{\prime\prime}(t) \ge 0\ \forall t \in I$
\end{theorem}

\begin{proof}
  Оба направления идут через промежуточное условие: $\phi^{\prime}$ не убывает на $I$. Для дифференцируемой $\phi^{\prime}$ это равносильно $\phi^{\prime\prime} \ge 0$ (признак монотонности). Поэтому достаточно доказать: $\phi$~--- выпуклая $\Leftrightarrow$ $\phi^{\prime}$ не убывает.

  \proofpart{$\Rightarrow$}
  Пусть $\phi$~--- выпуклая, $s, t \in I$. Запишем неравенство из \hyperref[note:2F475EE4-2630-45FD-8DD2-8EA9D48059E4]{одномерного дифференциального критерия} дважды:
  \[
  \phi(t) \ge \phi(s) + \phi^{\prime}(s)(t-s)
  \]
  \[
  \phi(s) \ge \phi(t) + \phi^{\prime}(t)(s-t)
  \]
  Сложим их и сократим $\phi(s) + \phi(t)$:
  \[
  0 \ge \phi^{\prime}(s)(t-s) + \phi^{\prime}(t)(s-t) = -(\phi^{\prime}(t) - \phi^{\prime}(s))(t-s)
  \]
  То есть $(\phi^{\prime}(t) - \phi^{\prime}(s))(t-s) \ge 0$: при $t > s$ получаем $\phi^{\prime}(t) \ge \phi^{\prime}(s)$~--- производная не убывает.

  \proofpart{$\Leftarrow$}
  Пусть $\phi^{\prime}$ не убывает, $s, t \in I$, $s \ne t$. По теореме Лагранжа существует $\xi$ между $s$ и $t$:
  \[
  \phi(t) - \phi(s) = \phi^{\prime}(\xi)(t-s)
  \]
  Если $t > s$, то $\xi > s \Rightarrow \phi^{\prime}(\xi) \ge \phi^{\prime}(s)$, и, умножая на $t - s > 0$, получаем $\phi^{\prime}(\xi)(t-s) \ge \phi^{\prime}(s)(t-s)$.
  Если $t < s$, то $\xi < s \Rightarrow \phi^{\prime}(\xi) \le \phi^{\prime}(s)$, и, умножая на $t - s < 0$, снова получаем $\phi^{\prime}(\xi)(t-s) \ge \phi^{\prime}(s)(t-s)$.
  В обоих случаях
  \[
  \phi(t) \ge \phi(s) + \phi^{\prime}(s)(t-s)
  \]
  а это по \hyperref[note:2F475EE4-2630-45FD-8DD2-8EA9D48059E4]{одномерному дифференциальному критерию} означает выпуклость $\phi$.
\end{proof}
  % Критерий выпуклости функции одной переменной через вторую производную -- одномерные случаи, к которым сводятся теоремы 1.2, 1.3

\InputIfFileExists{text/1E03F05F-22CC-4D55-BCF1-EBB5774765F0-before}{}{}
% generated by roam.py handbook export -- не редактировать, править заметки
% 20260921124320-теорема_ферма_необходимое_условие_экстремума.org (body ced7b28cf760)

\begin{theorem}[{Теорема Ферма (необходимое условие экстремума)}]\label[theorem]{note:1E03F05F-22CC-4D55-BCF1-EBB5774765F0}
  Пусть $g: (a;b) \rightarrow \R$, $t_0 \in (a;b)$~--- точка локального экстремума $g$, и $g$ дифференцируема в точке $t_0$.
  Тогда $g^{\prime}(t_0) = 0$
\end{theorem}

\begin{remark}
  Важно, что $t_0$~--- \emph{внутренняя} точка: к ней можно подойти с обеих сторон. На конце отрезка теорема неверна: $g(t) = t$ на $[0;1]$ имеет минимум в $t_0 = 0$, но $g^{\prime}(0) = 1$.
\end{remark}

\begin{remark}
  Обратное неверно: у $g(t) = t^3$ производная $g^{\prime}(0) = 0$, но экстремума в нуле нет. Условие $g^{\prime}(t_0) = 0$ только \emph{необходимое}.
\end{remark}

\begin{proof}
  Пусть $t_0$~--- точка локального минимума (для максимума применим то же рассуждение к $-g$).
  Тогда $\exists \delta > 0$: $g(t) \ge g(t_0)$ при $\abs{t - t_0} < \delta$.

  Посмотрим на разностное отношение с двух сторон от $t_0$. Числитель $g(t) - g(t_0) \ge 0$, а знак знаменателя зависит от стороны:
  \[
  t > t_0:\quad \frac{g(t) - g(t_0)}{t - t_0} \ge 0, \qquad t < t_0:\quad \frac{g(t) - g(t_0)}{t - t_0} \le 0
  \]
  Производная в $t_0$ существует, значит оба односторонних предела разностного отношения существуют и равны $g^{\prime}(t_0)$. Нестрогие неравенства сохраняются в пределе:
  \[
  g^{\prime}(t_0) = \lim_{t \rightarrow t_0 + 0} \frac{g(t) - g(t_0)}{t - t_0} \ge 0, \qquad
  g^{\prime}(t_0) = \lim_{t \rightarrow t_0 - 0} \frac{g(t) - g(t_0)}{t - t_0} \le 0
  \]
  Следовательно $g^{\prime}(t_0) = 0$.
\end{proof}
  % Теорема Ферма (необходимое условие экстремума) -- для $\Leftarrow$ в следствии 1.1

\InputIfFileExists{text/771B93FA-654D-4233-9487-B4FDD356DDFE-before}{}{}
% generated by roam.py handbook export -- не редактировать, править заметки
% 20260915184447-положительно_определенная_симметричная_матрица.org (body cde717804ad7)

\begin{definition}[{Положительно определенная симметричная матрица}]\label[definition]{note:771B93FA-654D-4233-9487-B4FDD356DDFE}
  Если матрица $A$~--- симметрична и $\forall x \ne 0\ x^TAx > 0$, то матрица положительно определена.
\end{definition}

\begin{remark}
  Если неравенство нестрогое, $\forall x\ x^TAx \ge 0$, то матрица \emph{неотрицательно определена} (пишут $A \ge 0$). Все свойства те же, но со сменой строгости: например, \hyperref[note:025A8419-962C-4892-A4F8-FB496AAA938C]{собственные значения} неотрицательны, а не положительны.
\end{remark}
  % Положительно определенная симметричная матрица

\InputIfFileExists{text/025A8419-962C-4892-A4F8-FB496AAA938C-before}{}{}
% generated by roam.py handbook export -- не редактировать, править заметки
% 20260915183647-собственные_значения_положительно_определенных_симметричных_матриц.org (body 32f3a0ed8262)

\begin{theorem}[{Собственные значения положительно определенных, симметричных матриц}]\label[theorem]{note:025A8419-962C-4892-A4F8-FB496AAA938C}
  У \hyperref[note:771B93FA-654D-4233-9487-B4FDD356DDFE]{положительно определенной симметричной матрицы} все ее \hyperref[note:9D2905FC-8E68-4E6F-9243-364269F795B0]{собственные значения} являются положительными вещественными числами.
\end{theorem}

\begin{proof}
  Пусть $A$~--- положительно определенная симметричная матрица, $\lambda$~--- ее собственное значение. По \hyperref[note:4F4E5575-ABE6-4D9B-9C61-B76A3AA9701C]{теореме о симметричных матрицах} $\lambda \in \R$.

  Собственный вектор можно взять вещественным: матрица $A - \lambda I$~--- вещественная и вырожденная, значит у системы $(A - \lambda I) q = 0$ есть ненулевое решение $q \in \R^n$.

  Подставим $x = q$ в определение положительной определенности:
  \[
  0 < q^T A q = q^T (\lambda q) = \lambda q^T q = \lambda \norm{q}^2
  \]
  Поделим на $\norm{q}^2 > 0$ и получим $\lambda > 0$.
\end{proof}
  % Собственные значения положительно определенных, симметричных матриц -- смысл записи $\nabla^2 f \ge 0$ в теореме 1.3

\InputIfFileExists{text/9D2905FC-8E68-4E6F-9243-364269F795B0-before}{}{}
% generated by roam.py handbook export -- не редактировать, править заметки
% 20260915181503-собственное_значение.org (body fa969fdb0ac3)

\begin{definition}[{Собственное значение}]\label[definition]{note:9D2905FC-8E68-4E6F-9243-364269F795B0}
  Число $\lambda \in \C$ является \textbf{собственным значением} \emph{квадратной} матрицы $A$, если существует такой \emph{вектор} $q$, что $A q = \lambda q$ и $q$~--- \emph{не нулевой}

  Такой вектор называется \hyperref[note:57D369D6-D748-41A2-9BDA-38D2FFC6601A]{собственным вектором} матрицы.
\end{definition}
  % Собственное значение

\InputIfFileExists{text/4F4E5575-ABE6-4D9B-9C61-B76A3AA9701C-before}{}{}
% generated by roam.py handbook export -- не редактировать, править заметки
% 20260915183320-собственные_значения_симметричных_матриц.org (body b0b3509b735f)

\begin{theorem}[{Собственные значения симметричных матриц}]\label[theorem]{note:4F4E5575-ABE6-4D9B-9C61-B76A3AA9701C}
  У вещественных симметричных матриц все \hyperref[note:9D2905FC-8E68-4E6F-9243-364269F795B0]{собственные значения}~--- вещественны
\end{theorem}

\begin{proof}
  Пусть $A = A^T$~--- вещественная матрица, $Aq = \lambda q$, $q \ne 0$. Здесь $\lambda \in \C$ и $q \in \C^n$.
  Обозначим $q^* = \overline{q}^T$ и посчитаем число $q^* A q$:
  \[
  q^* A q = q^* (\lambda q) = \lambda q^* q = \lambda \norm{q}^2
  \]
  Это число совпадает со своим сопряженным:
  \[
  \overline{q^* A q} = (q^* A q)^* = q^* A^* q = q^* A q
  \]
  так как $A^* = \overline{A}^T = A^T = A$ ($A$~--- вещественная и симметричная). Значит $q^* A q \in \R$.

  Так как $\norm{q}^2 > 0$~--- вещественное число, то
  \[
  \lambda = \frac{q^* A q}{\norm{q}^2} \in \R
  \]
\end{proof}
  % Собственные значения симметричных матриц

\InputIfFileExists{text/57D369D6-D748-41A2-9BDA-38D2FFC6601A-before}{}{}
% generated by roam.py handbook export -- не редактировать, править заметки
% 20260915181913-собственныи_вектор.org (body fbd08d1eac3a)

\begin{definition}[{Собственный вектор}]\label[definition]{note:57D369D6-D748-41A2-9BDA-38D2FFC6601A}
  $q$~--- собственный вектор матрицы $A$, если он \emph{не нулевой} и существует такое \hyperref[note:9D2905FC-8E68-4E6F-9243-364269F795B0]{Собственное значение} $\lambda$, что $Aq = \lambda q$
\end{definition}
  % Собственный вектор

\section{К лекции 2}\label{aux:2}

\InputIfFileExists{text/957FA09D-8883-40A5-B4C7-237D9F815A42-before}{}{}
% generated by roam.py handbook export -- не редактировать, править заметки
% 20260918141614-дифференцируемая_функция.org (body 2150a53e17b6)

\begin{definition}[{Дифференцируемая функция}]\label[definition]{note:957FA09D-8883-40A5-B4C7-237D9F815A42}
  Пусть $E \subset \R^n$~--- открытое множество, $f: E \rightarrow \R^m$, $x \in E$.
  $f$ \textbf{дифференцируема} в точке $x$, если существует линейное отображение (матрица $A$ размера $m \times n$) такое, что
  \[
  f(x + h) = f(x) + A h + o(\norm{h}),\ h \rightarrow 0
  \]
  Матрица $A$ единственна, называется \emph{производной} (матрицей Якоби) $f$ в точке $x$ и обозначается $f^{\prime}(x)$; ее элементы~--- частные производные: $A_{ij} = \frac{\partial f_i}{\partial x_j}(x)$.
  $f$ дифференцируема на $E$, если она дифференцируема в каждой точке $E$.

  \begin{itemize}
    \item При $m = 1$ производная~--- строка $1 \times n$, а ее транспонированная запись~--- \hyperref[note:1756907B-1FF2-4F62-8553-759B8509281C]{градиент}: $f^{\prime}(x) h = \left\langle \nabla f(x); h \right\rangle$.
    \item Для $g = \nabla f: E \rightarrow \R^n$ производная $g^{\prime}(x)$~--- это \hyperref[note:592AD816-1D80-48E6-8399-8524FBC1E8F3]{гессиан} $\nabla^2 f(x)$.
    \item $f \in C^k(E)$~--- $f$ $k$ раз \emph{непрерывно} дифференцируема: все частные производные до порядка $k$ существуют и непрерывны. «Дважды дифференцируема» в заметках курса обычно означает $f \in C^2$.
  \end{itemize}
\end{definition}

\begin{remark}
  Существование частных производных не влечет дифференцируемость: функция
  \[
  f(x_1, x_2) = \frac{x_1 x_2}{x_1^2 + x_2^2},\quad f(0) = 0
  \]
  имеет в нуле обе частные производные (равные $0$), но даже не непрерывна там (на прямой $x_1 = x_2$ она равна $\frac{1}{2}$). Достаточное условие: если частные производные существуют в окрестности $x$ и непрерывны в $x$, то $f$ дифференцируема в $x$; в частности, $C^1 \Rightarrow$ дифференцируема.
\end{remark}

\begin{proposition}[{Свойство: дифференцируемая функция непрерывна}]\label[proposition]{note:957FA09D-8883-40A5-B4C7-237D9F815A42:1}
  Из определения
  \[
  \norm{f(x + h) - f(x)} \le \norm{A h} + \norm{o(\norm{h})} \le \norm{A} \cdot \norm{h} + o(\norm{h}) \rightarrow 0,\ h \rightarrow 0
  \]
  где $\norm{A}$~--- операторная норма. Значит, $f(x + h) \rightarrow f(x)$.
\end{proposition}

\paragraph*{Единственность производной}
Пусть $A$ и $B$ обе подходят. Вычитая разложения, получаем $(A - B) h = o(\norm{h})$. Возьмем $h = t e$ для фиксированного $e$ и $t \rightarrow 0+$: $(A - B) e = \frac{o(t)}{t} \rightarrow 0$, то есть $(A - B) e = 0$ для любого $e$, и $A = B$.
  % Дифференцируемая функция

\InputIfFileExists{text/3BE0BBCE-AAB0-4CB4-8FA6-45277C0889CF-before}{}{}
% generated by roam.py handbook export -- не редактировать, править заметки
% 20260923124012-спектральная_теорема_для_симметричных_матриц.org (body 3c7c31069135)

\begin{theorem}[{Спектральная теорема для симметричных матриц}]\label[theorem]{note:3BE0BBCE-AAB0-4CB4-8FA6-45277C0889CF}
  Пусть $A \in \R^{n \times n}$~--- симметричная матрица ($A = A^T$).
  Тогда в $\R^n$ существует ортонормированный базис из \hyperref[note:57D369D6-D748-41A2-9BDA-38D2FFC6601A]{собственных векторов} $A$:
  \[
  A q_i = \lambda_i q_i,\quad \left\langle q_i; q_j \right\rangle = \begin{cases} 1, & i = j \\ 0, & i \ne j \end{cases},\quad \lambda_i \in \R
  \]
\end{theorem}

\begin{remark}
  Матричная форма: $A = Q \Lambda Q^T$, где $Q = (q_1 \mid \dots \mid q_n)$~--- ортогональная матрица ($Q^T Q = I$), $\Lambda = \mathrm{diag}(\lambda_1, \dots, \lambda_n)$.
\end{remark}

\begin{remark}
  Главное следствие для оптимизации: в этом базисе квадратичная форма становится суммой квадратов. Если $h = \sum\limits_{i=1}^n c_i q_i$, то $\norm{h}^2 = \sum\limits_{i=1}^n c_i^2$ и
  \[
  \left\langle Ah; h \right\rangle = \sum\limits_{i=1}^n \lambda_i c_i^2, \qquad \lambda_{\min} \norm{h}^2 \le \left\langle Ah; h \right\rangle \le \lambda_{\max} \norm{h}^2
  \]
  Отсюда $A \ge \mu I \Leftrightarrow \lambda_{\min}(A) \ge \mu$ (см. \hyperref[note:94DADA28-87E4-4E36-B289-8415BCE750FC]{критерий сильной выпуклости в терминах вторых производных}).
\end{remark}

\begin{proof}
  Индукция по $n$. База $n = 1$: подходит $q_1 = 1$.

  \proofpart{Первый собственный вектор}
  По теореме о \hyperref[note:4F4E5575-ABE6-4D9B-9C61-B76A3AA9701C]{собственных значениях симметричных матриц} все \hyperref[note:9D2905FC-8E68-4E6F-9243-364269F795B0]{собственные значения} $A$ вещественны. Возьмем одно из них, $\lambda_1$.
  Матрица $A - \lambda_1 I$ вещественная и вырожденная, поэтому у нее есть ненулевое вещественное решение $(A - \lambda_1 I) q = 0$. Положим $q_1 = \frac{q}{\norm{q}}$.

  \proofpart{Ортогональное дополнение инвариантно}
  Пусть $W = \left\{ h \in \R^n : \left\langle h; q_1 \right\rangle = 0 \right\}$, $\dim W = n - 1$. Для $h \in W$ по симметричности $A$:
  \[
  \left\langle Ah; q_1 \right\rangle = h^T A^T q_1 = h^T A q_1 = \lambda_1 \left\langle h; q_1 \right\rangle = 0
  \]
  то есть $A W \subset W$.

  \proofpart{Переход}
  Выберем в $W$ ортонормированный базис $u_1, \dots, u_{n-1}$ и составим из него матрицу $U = (u_1 \mid \dots \mid u_{n-1})$ размера $n \times (n-1)$, $U^T U = I$.
  Для $w \in W$ верно $w = U U^T w$, а так как $A U$ состоит из векторов $W$, то $A U = U B$, где $B = U^T A U$.
  Матрица $B$ размера $(n-1) \times (n-1)$ симметрична: $B^T = U^T A^T U = B$. По предположению индукции у нее есть ортонормированный базис из собственных векторов $p_2, \dots, p_n$, $B p_i = \lambda_i p_i$.

  Положим $q_i = U p_i$ при $i = 2, \dots, n$. Тогда
  \[
  A q_i = A U p_i = U B p_i = \lambda_i U p_i = \lambda_i q_i, \qquad \left\langle q_i; q_j \right\rangle = p_i^T U^T U p_j = \left\langle p_i; p_j \right\rangle
  \]
  то есть $q_2, \dots, q_n$~--- ортонормированные собственные векторы $A$, лежащие в $W$, и потому ортогональные $q_1$.
  Вместе с $q_1$ получаем $n$ ортонормированных векторов~--- базис $\R^n$.
\end{proof}
  % Спектральная теорема для симметричных матриц

\InputIfFileExists{text/D0601ACD-8C91-4A93-B0E4-AC0D529C6A5C-before}{}{}
% generated by roam.py handbook export -- не редактировать, править заметки
% 20260923123518-неравенство_коши_буняковского_шварца.org (body 56b35acc942e)

\begin{theorem}[{Неравенство Коши-Буняковского-Шварца}]\label[theorem]{note:D0601ACD-8C91-4A93-B0E4-AC0D529C6A5C}
  Пусть $X$~--- \hyperref[note:1B626AA4-35C8-4C7C-9D3D-D8C26D98F40C]{векторное пространство} над $\R$ со \hyperref[note:F2942E98-6115-4600-AA3F-E0F687D9C732]{скалярным произведением}, $\norm{x} = \sqrt{\left\langle x; x \right\rangle}$.
  Тогда
  \[
  \abs{\left\langle x; y \right\rangle} \le \norm{x} \norm{y}\quad \forall x, y \in X
  \]
  причем равенство достигается тогда и только тогда, когда $x$ и $y$ линейно зависимы.
\end{theorem}

\begin{remark}
  Раскрывая модуль, получаем двустороннюю оценку $-\norm{x} \norm{y} \le \left\langle x; y \right\rangle \le \norm{x} \norm{y}$. В оптимизации чаще нужна левая часть: линейное слагаемое $\left\langle g; y \right\rangle$ не может уменьшить значение сильнее, чем на $\norm{g} \norm{y}$.
\end{remark}

\begin{proof}
  Если $y = 0$, обе части равны нулю. Пусть $y \ne 0$. Для любого $t \in \R$ по неотрицательности скалярного квадрата
  \[
  0 \le \left\langle x - ty; x - ty \right\rangle = \norm{x}^2 - 2t \left\langle x; y \right\rangle + t^2 \norm{y}^2
  \]
  Справа~--- квадратный трехчлен по $t$ с положительным старшим коэффициентом $\norm{y}^2$, который не принимает отрицательных значений. Значит, его дискриминант неположителен:
  \[
  4 \left\langle x; y \right\rangle^2 - 4 \norm{x}^2 \norm{y}^2 \le 0
  \]
  откуда $\abs{\left\langle x; y \right\rangle} \le \norm{x} \norm{y}$.

  \proofpart{Случай равенства}
  Равенство означает, что дискриминант равен нулю, то есть у трехчлена есть корень $t_0$: $\left\langle x - t_0 y; x - t_0 y \right\rangle = 0$, а значит $x = t_0 y$.
  Обратно, при $x = t_0 y$ обе части равны $\abs{t_0} \norm{y}^2$.
\end{proof}
  % Неравенство Коши-Буняковского-Шварца -- теоремы 2.3, 2.4

\InputIfFileExists{text/AD362344-ECDC-4627-B344-8356744E6E53-before}{}{}
% generated by roam.py handbook export -- не редактировать, править заметки
% 20260923142500-формула_ньютона_леибница.org (body 59c10fef8a22)

\begin{theorem}[{Формула Ньютона-Лейбница}]\label[theorem]{note:AD362344-ECDC-4627-B344-8356744E6E53}
  Пусть $g: [a;b] \rightarrow \R$~--- непрерывная функция, $G$~--- ее первообразная на $[a;b]$ ($G^{\prime} = g$). Тогда
  \[
  \int\limits_a^b g(t)\, dt = G(b) - G(a)
  \]
  В заметках курса чаще нужна форма для $\phi \in C^1[a;b]$ (первообразная для $\phi^{\prime}$~--- сама $\phi$):
  \[
  \phi(b) - \phi(a) = \int\limits_a^b \phi^{\prime}(t)\, dt
  \]
\end{theorem}

\begin{remark}
  Для вектор-функции $\phi: [a;b] \rightarrow \R^m$ формула верна покоординатно, интеграл понимается как вектор интегралов координат.
\end{remark}

\begin{remark}
  Типичное применение: для $\phi(t) = f(x + t(y - x))$ получаем
  \[
  f(y) - f(x) = \int\limits_0^1 \left\langle \nabla f(x + t(y-x)); y - x \right\rangle dt
  \]
  так доказываются \hyperref[note:379E4CAD-A44B-4631-8DAF-F3DC708EAEDE]{усиленное неравенство для функций с липшицевым градиентом} и \hyperref[note:81C8B68E-DAAD-4B86-BAEB-AB0E1C07442F]{аппроксимация функции с липшицевым гессианом} (там формула применяется дважды, поэтому нужна непрерывность $\phi^{\prime\prime}$, то есть $f \in C^2$).
\end{remark}

\begin{proof}
  \proofpart{Интеграл с переменным верхним пределом~--- первообразная}
  Пусть $F(x) = \int\limits_a^x g(t)\, dt$. Для $x \in [a;b]$ и малого $h \ne 0$
  \[
  \frac{F(x + h) - F(x)}{h} - g(x) = \frac{1}{h} \int\limits_x^{x+h} \big( g(t) - g(x) \big) dt
  \]
  поэтому
  \[
  \abs{\frac{F(x + h) - F(x)}{h} - g(x)} \le \sup\limits_{\abs{t - x} \le \abs{h}} \abs{g(t) - g(x)} \rightarrow 0,\ h \rightarrow 0
  \]
  по непрерывности $g$ в точке $x$. Значит, $F^{\prime}(x) = g(x)$.

  \proofpart{Две первообразные отличаются на константу}
  $(G - F)^{\prime} = g - g = 0$ на $[a;b]$, поэтому по теореме Лагранжа $G - F \equiv C$. Тогда
  \[
  G(b) - G(a) = F(b) - F(a) = \int\limits_a^b g(t)\, dt - 0
  \]
\end{proof}
  % Формула Ньютона-Лейбница -- теоремы 2.3, 2.4

\InputIfFileExists{text/5821D7FE-1F9F-4578-BA2B-D60FDB0C69D3-before}{}{}
% generated by roam.py handbook export -- не редактировать, править заметки
% 20260923130752-теорема_лагранжа_для_векторнозначных_функции.org (body 0b41906a3917)

\begin{theorem}[{Лагранжа для векторнозначных функций}]\label[theorem]{note:5821D7FE-1F9F-4578-BA2B-D60FDB0C69D3}
  Пусть $E \subset \R^n$~--- открытое \hyperref[note:316B5BFA-8B9F-4267-9086-D6593AAB5F9B]{выпуклое множество}, $g: E \rightarrow \R^m$~--- \hyperref[note:957FA09D-8883-40A5-B4C7-237D9F815A42]{дифференцируемое} отображение.
  Тогда
  \[
  \forall a, b \in E\ \exists c \in [a;b]:\quad \norm{g(b) - g(a)} \le \norm{g^{\prime}(c)} \cdot \norm{b - a}
  \]
  где $g^{\prime}(c)$~--- матрица Якоби $m \times n$, а $\norm{g^{\prime}(c)}$~--- ее спектральная (операторная) норма: $\norm{A} = \sup\limits_{h \ne 0} \frac{\norm{Ah}}{\norm{h}}$.
\end{theorem}

\begin{remark}
  В отличие от одномерного случая, \emph{равенства} $g(b) - g(a) = g^{\prime}(c)(b - a)$ может не быть ни при каком $c$. Например, $g(t) = (\cos t, \sin t)$ на $[0; 2\pi]$: $g(2\pi) - g(0) = 0$, но $g^{\prime}(c) = (-\sin c, \cos c) \ne 0$ при всех $c$.
\end{remark}

\begin{remark}
  Для $g = \nabla f$ матрица Якоби~--- это \hyperref[note:592AD816-1D80-48E6-8399-8524FBC1E8F3]{гессиан} $\nabla^2 f$.
\end{remark}

\begin{proof}
  Пусть $v = g(b) - g(a)$. Если $v = 0$, доказывать нечего. Иначе сведем все к скалярной функции
  \[
  \phi(t) = \left\langle v; g(a + t(b - a)) \right\rangle, \quad t \in [0;1]
  \]
  Она дифференцируема и $\phi^{\prime}(t) = \left\langle v; g^{\prime}(a + t(b-a))(b - a) \right\rangle$.
  По одномерной теореме Лагранжа $\exists \theta \in (0;1)$: $\phi(1) - \phi(0) = \phi^{\prime}(\theta)$. Положим $c = a + \theta(b - a) \in [a;b]$.

  Слева $\phi(1) - \phi(0) = \left\langle v; g(b) - g(a) \right\rangle = \norm{v}^2$. Справа по \hyperref[note:D0601ACD-8C91-4A93-B0E4-AC0D529C6A5C]{КБШ} и определению операторной нормы
  \[
  \left\langle v; g^{\prime}(c)(b - a) \right\rangle \le \norm{v} \cdot \norm{g^{\prime}(c)(b - a)} \le \norm{v} \cdot \norm{g^{\prime}(c)} \cdot \norm{b - a}
  \]
  Итого $\norm{v}^2 \le \norm{v} \cdot \norm{g^{\prime}(c)} \cdot \norm{b - a}$; делим на $\norm{v} > 0$.
\end{proof}
  % Теорема Лагранжа для векторнозначных функций -- следствие 2.1

\InputIfFileExists{text/82C70263-336A-49F5-82DB-9D25C84D6A6D-before}{}{}
% generated by roam.py handbook export -- не редактировать, править заметки
% 20260923123518-теорема_веиерштрасса_о_достижении_экстремума.org (body fdabd43a4994)

\begin{theorem}[{Вейерштрасса о достижении экстремума}]\label[theorem]{note:82C70263-336A-49F5-82DB-9D25C84D6A6D}
  Пусть $K \subset \R^n$~--- непустое компактное (замкнутое и ограниченное) множество, $f: K \rightarrow \R$~--- непрерывная функция.
  Тогда $f$ ограничена на $K$ и достигает на нем наименьшего и наибольшего значений:
  \[
  \exists x_{*}, x^{*} \in K:\quad f(x_{*}) \le f(x) \le f(x^{*})\quad \forall x \in K
  \]

  \begin{itemize}
    \item $K$ не замкнуто: $f(x) = x$ на $(0;1]$~--- $\inf f = 0$ не достигается;
    \item $K$ не ограничено: $f(x) = e^{-x}$ на $[0;+\infty)$~--- $\inf f = 0$ не достигается;
    \item $f$ разрывна: $f(x) = x$ при $x \in [0;1)$, $f(1) = 0$~--- $\sup f = 1$ не достигается.
  \end{itemize}
\end{theorem}

\begin{remark}
  Все условия существенны:
\end{remark}

\begin{remark}
  Типичное применение в оптимизации: если $f$ растет на бесконечности, то вне большого шара ее значения заведомо больше, чем в какой-то точке внутри, и минимум достаточно искать на шаре (см. \hyperref[note:51CC880D-BF19-43E2-9599-FB31A9A7367F]{единственность минимума сильно выпуклой функции}).
\end{remark}

\begin{proof}
  Докажем для минимума; для максимума применим то же рассуждение к $-f$.

  \proofpart{Ограниченность снизу}
  Пусть $f$ не ограничена снизу: тогда $\forall k \in \mathbb{N}\ \exists x_k \in K:\ f(x_k) < -k$.
  $K$ ограничено, поэтому по теореме Больцано~--- Вейерштрасса у $(x_k)$ есть сходящаяся подпоследовательность $x_{k_j} \rightarrow \bar{x}$, а так как $K$ замкнуто, $\bar{x} \in K$.
  По непрерывности $f(x_{k_j}) \rightarrow f(\bar{x})$~--- конечное число, но $f(x_{k_j}) < -k_j \rightarrow -\infty$. Противоречие.

  \proofpart{Достижение}
  Пусть $m = \inf_K f > -\infty$. По определению инфимума $\forall k\ \exists x_k \in K:\ m \le f(x_k) < m + \frac{1}{k}$.
  Как и выше, выделим подпоследовательность $x_{k_j} \rightarrow x_{*} \in K$. По непрерывности
  \[
  f(x_{*}) = \lim_{j \rightarrow \infty} f(x_{k_j}) = m
  \]
  то есть $x_{*}$~--- точка минимума $f$ на $K$.
\end{proof}
  % Теорема Вейерштрасса о достижении экстремума -- единственность минимума
